\chapter*{Appendix I: Key Risk Indicator Design — Reference Framework}
\label{chap:Appendix I: Key Risk Indicator Design — Reference Framework}

\paragraph{A reference framework for designing Key Risk Indicators, drawing on the principles discussed in Chapter~\ref{chap:Step 5 - Recording, Reviewing, and Monitoring}.}

\section*{The KRI Design Checklist}

\paragraph{For each proposed KRI, the following questions should be answered:}

\begin{itemize}
    \bitem{Relevance}{does this metric directly reflect the risk it is designed to monitor?}
    \bitem{Leading or coincident}{does the metric change in advance of the risk materialising (leading) or at the same time (coincident)? Leading indicators are preferred.}
    \bitem{Data availability}{is the data required to calculate this metric available, reliable, and accessible on the required frequency?}
    \bitem{Ownership}{is there a named individual responsible for monitoring this KRI and acting on threshold breaches?}
    \bitem{Thresholds}{are clear threshold levels defined — the level at which the indicator moves from green to amber, and from amber to red?}
    \bitem{Escalation}{is the escalation pathway for amber and red status clearly defined?}
    \bitem{Review}{is the KRI reviewed periodically to confirm it remains the right metric for the risk?}
\end{itemize}

\begin{table}[H]
  \centering
  \renewcommand{\arraystretch}{1.5}
  \begin{tabularx}{\textwidth}{ | >{\raggedright\arraybackslash}p{0.35\textwidth} 
                                | >{\raggedright\arraybackslash}X | }
    \hline
    \textbf{Risk Category} & \textbf{Example KRIs} \\
    \hline
    \textbf{Credit risk} & Portfolio delinquency rate; percentage of loans in watchlist; concentration by single borrower\\
    \hline
    \textbf{Liquidity risk} & Liquidity coverage ratio; short-term funding ratio; counterparty concentration in funding\\
    \hline
    \textbf{Operational risk} & System availability statistics; processing error rates; incident frequency by category\\
    \hline
    \textbf{Conduct risk} & Customer complaint volumes and trends; upheld complaint rate; mystery shopping scores\\
    \hline
    \textbf{Cybersecurity risk} & Phishing simulation failure rate; unpatched critical vulnerabilities; mean time to detect incidents\\
    \hline
    \textbf{People risk} & Staff turnover rate; vacancy rate in key roles; whistleblowing report volumes\\
    \hline
    \textbf{Third-party risk} & SLA breach frequency; supplier financial health scores; overdue third-party assessments\\
    \hline
    \textbf{Regulatory risk} & Regulatory correspondence frequency; overdue regulatory submissions; internal audit findings\\
    \hline
    \textbf{Reputational risk} & Net Promoter Score; social media sentiment score; adverse media mentions\\
    \hline
    \textbf{Climate risk} & Scope 1, 2, and 3 emissions trajectory; physical risk exposure score; transition risk score\\
    \hline
  \end{tabularx}
  \caption*{KRI Categories — Examples by Risk Type}
\end{table}
